Los experimentos fueron ejecutados en un entorno Linux
sobre WSL (Windows Subsystem for Linux), utilizando un
procesador AMD Ryzen 3 7320U with Radeon Graphics (2.40 GHz),
Sistema operativo de 64 bits, procesador x64 y 8 GB de memoria
RAM. Los programas principales fueron implementados en C++17
y compilados con \texttt{g++} utilizando el nivel de
optimización predeterminado vía \texttt{Makefile}. 

\vspace{0.4cm}

La medición de tiempo se realizó con la librería estándar
\texttt{<chrono>} (\texttt{high\_resolution\_clock}),
mientras que el consumo de memoria \textit{Peak RSS} se
registró mediante la función \texttt{getrusage()} de la
cabecera POSIX \texttt{<sys/resource.h>}. 

\vspace{0.4cm}

Para las pruebas de ordenamiento se generaron arreglos
unidimensionales mediante un script de python con tamaños 
$n \in \{10^1, 10^3, 10^5, 10^7\}$, considerando tres tipos
de orden inicial (\textit{ascendente},
\textit{descendente}, \textit{aleatorio}) y dos dominios de
valores ($D1: [0, 9]$ y $D7: [0, 10^7]$),
con lo cual se construye el dataset de arreglos. 

\vspace{0.4cm}

Para las pruebas de multiplicación de matrices, se generaron 
matrices mediante un script de python utilizando dimensiones
$n \in \{2^4, 2^6, 2^8, 2^{10}\}$ en tres estructuras
(\textit{densa}, \textit{diagonal}, \textit{dispersa})
sobre los dominios $D0: \{0, 1\}$ y $D10: [0, 9]$, con lo cual
se construye el dataset de matrices. 

\vspace{0.4cm}

Las visualizaciones graficas fueron generadas mediante
scripts de Python usando las libreria \texttt{pandas} y
\texttt{seaborn}.